% ============================================================
% Mainline summary: Flow-GRPO credit assignment
% ============================================================

\section{主线总结：Flow-GRPO 信用分配四篇阶段总结}

\begin{conceptbox}
\textbf{覆盖论文：} DenseGRPO、Chunk-GRPO、TempFlow-GRPO、PCPO。\\
\textbf{总结日期：} 2026-04-28。\\
\textbf{当前判断：} 这四篇构成一个比较清晰的主线：\textbf{Flow-GRPO 后训练的 credit assignment 不能只看 terminal reward，也不能只看 advantage；policy ratio / logprob scale 本身也会改变每个 step 的实际更新强度。}
\end{conceptbox}

\subsection{0. 最核心的一句话}

标准 Flow-GRPO 的问题可以压缩成一个式子：

\[
\text{实际每个 step 的更新强度}
\approx
\underbrace{\text{advantage / credit}}_{\text{reward side}}
\times
\underbrace{\text{policy ratio / logprob scale}}_{\text{policy side}}.
\]

更具体地，PPO / GRPO 的 step-level objective 里有：

\[
\min
\left(
\rho_t^i A_i,\,
\mathrm{clip}(\rho_t^i,1-\epsilon,1+\epsilon)A_i
\right).
\]

所以不能只问：

\[
A_i \text{ 应该怎么分给每个 timestep？}
\]

还要问：

\[
\rho_t^i \text{ 本身是不是已经在放大或压小某些 timestep？}
\]

\begin{summarybox}
\textbf{阶段性结论：} 好的 Flow-GRPO 后训练，需要同时校准 \textbf{advantage side} 和 \textbf{ratio side}。前者决定“奖励应该分给谁”，后者决定“同样的奖励乘到 objective 里以后，实际更新强度会不会偏”。
\end{summarybox}

\subsection{1. 第一类问题：advantage 分布不平衡}

标准 Flow-GRPO / DanceGRPO 的做法是：同一个 prompt 采样 \(G\) 条 trajectories，最终每条 trajectory 得到一个 terminal reward，再归一化成 trajectory-level advantage：

\[
A_i
=
\frac{
R(x_0^i,c)-\mathrm{mean}_{j}R(x_0^j,c)
}{
\mathrm{std}_{j}R(x_0^j,c)+\epsilon
}.
\]

然后每个 denoising step 都拿同一个 \(A_i\)：

\[
\sum_{t=1}^{T}
\min
\left(
\rho_t^i A_i,\,
\mathrm{clip}(\rho_t^i,1-\epsilon,1+\epsilon)A_i
\right).
\]

这等价于默认：

\[
\text{一条 trajectory 最终好}
\Rightarrow
\text{这条 trajectory 的每个 step 都好。}
\]

这个假设太强。因为最终 reward 高，不代表所有中间 steps 都贡献了正收益；最终 reward 低，也不代表所有中间 steps 都应该被惩罚。

\begin{warnbox}
\textbf{这就是 advantage 分布不平衡：} terminal advantage 被全局广播到所有 timesteps，导致真正有贡献的 step 和没贡献的 step 拿到同一个训练信号。它不是 reward model 不准的问题，而是 reward credit 的时间分配太粗。
\end{warnbox}

\subsection{2. DenseGRPO：最直接地修 advantage side}

DenseGRPO 的核心想法是把 terminal reward 拆成 step-wise reward gain：

\[
\Delta R_t
=
R_{t-1}-R_t.
\]

这里的 \(R_t\) 不是直接给 noisy latent 打分，而是从中间 latent \(x_t\) 出发，用 ODE denoising readout 到 clean image，再 decode、再用 reward model 打分。

所以 DenseGRPO 本质上在问：

\[
\text{经过第 \(t\) 个 denoising transition 后，最终可读出的 reward 提升了多少？}
\]

这比 terminal advantage 更合理，因为它更接近：

\[
A_i
\quad \rightarrow \quad
A_t^i.
\]

\begin{summarybox}
\textbf{DenseGRPO 的价值：} 它是四篇里对 advantage side 修得最直接的一篇。它试图估计每个 step 的 actual improvement，而不是把最终 reward 平均分配给所有 steps。
\end{summarybox}

但 DenseGRPO 的问题也很清楚：

\[
T+(T-1)+\cdots+1
=
\frac{T(T+1)}{2}.
\]

如果每个中间 latent 都要完整 ODE readout 到 clean image，成本接近 \(O(T^2)\)。所以 DenseGRPO 的核心收益很硬，工程成本也很硬。

\subsection{3. Chunk-GRPO：不直接修 reward，而是修优化粒度}

Chunk-GRPO 的判断是：单个 step 太细，整条 trajectory 太粗，所以可以用 chunk 作为中间粒度。

它把 denoising trajectory 切成：

\[
\mathrm{ch}_1,\mathrm{ch}_2,\ldots,\mathrm{ch}_K.
\]

然后把 step-level ratio 合成 chunk-level ratio：

\[
r_j^i(\theta)
=
\left(
\prod_{t\in \mathrm{ch}_j}
\frac{
p_\theta(x_{t-1}^i|x_t^i,c)
}{
p_{\mathrm{old}}(x_{t-1}^i|x_t^i,c)
}
\right)^{1/cs_j}.
\]

其中 \(1/cs_j\) 是几何平均，防止 chunk 越长 ratio 尺度越容易爆。

Chunk-GRPO 没有像 DenseGRPO 那样真正计算 chunk reward improvement。它更多是在说：

\[
\text{更新单位不一定要是单 step，也不应该是整条 trajectory。}
\]

\begin{summarybox}
\textbf{Chunk-GRPO 的价值：} 它让 credit 分配更平滑，减少 step-level 噪声，也提供了一个很实用的工程方向：如果 DenseGRPO 的 step-wise readout 太贵，可以尝试 chunk-level improvement。
\end{summarybox}

它的局限是：如果仍然用同一个 terminal \(A_i\)，那它只是缓解了 credit assignment，没有彻底解决“哪个 chunk 真正带来 reward improvement”的问题。

\subsection{4. TempFlow-GRPO：用 branch 做局部反事实}

TempFlow-GRPO 试图回答的是另一个问题：

\[
\text{某个 timestep 的选择，是否真的导致了最终结果差异？}
\]

它用 trajectory branching 的方式，在某个时刻 \(t\) 产生分支，然后观察不同 branch 的后续结果。直觉上，这比单纯看一条 trajectory 更接近反事实：

\[
\text{如果在同一个前缀状态下，换一个 branch，最终 reward 会怎样？}
\]

所以 TempFlow-GRPO 的 credit localization 不是 DenseGRPO 那种 step-wise reward gain，而是 branch-level counterfactual。

\begin{summarybox}
\textbf{TempFlow-GRPO 的价值：} 它提醒我们，credit 不只是 timestep 的问题，也和局部探索有关。一个 step 的重要性，最好通过“同一位置不同 branch 的后果差异”来判断。
\end{summarybox}

但它也有两个问题：

\begin{enumerate}[leftmargin=1.5em]
  \item branch 怎么设计会影响 credit 结论。
  \item 分支越多，成本越高；分支太少，反事实又不够充分。
\end{enumerate}

\subsection{5. PCPO：ratio side 也会影响 credit}

前面三篇主要都在处理 advantage side：

\[
A_i
\rightarrow
A_t^i
\quad \text{or} \quad
A_{\mathrm{chunk}}^i
\quad \text{or} \quad
A_{\mathrm{branch}}^i.
\]

PCPO 的关键提醒是：即使 advantage 先不变，policy ratio 本身也不是中性的。

Flow-GRPO 里每一步都有：

\[
\rho_t^i(\theta)
=
\frac{
p_\theta(x_{t-1}^i|x_t^i,c)
}{
p_{\theta_{\mathrm{old}}}(x_{t-1}^i|x_t^i,c)
}.
\]

PCPO 展开 Flow SDE 的 log-ratio 后指出，不同 timestep 的 log-ratio scale 会受到 step size、SDE variance、timestep position 等因素影响。也就是说，即使用同一个 \(A_i\)，不同 timestep 的实际更新强度也可能不一样。

可以把它理解成：

\[
\rho_t^i A_i
\quad
\text{里面的}
\quad
\rho_t^i
\quad
\text{已经带了一个隐式 timestep weight。}
\]

\begin{warnbox}
\textbf{PCPO 的关键意义：} credit assignment 不只发生在 \(A\) 里，也发生在 \(\rho\) 里。只修 advantage side，不看 ratio side，可能仍然会让某些 steps 被 objective 数学形式过度放大或压小。
\end{warnbox}

所以 PCPO 把我们的主线从：

\[
\text{如何分配 reward？}
\]

推进到：

\[
\text{如何校准每个 step 的实际 policy update 强度？}
\]

\subsection{6. 四篇放在一起看}

\begin{center}
\resizebox{\linewidth}{!}{
\begin{tabular}{llll}
\toprule
\textbf{论文} & \textbf{修哪一侧} & \textbf{核心机制} & \textbf{主要局限}\\
\midrule
DenseGRPO & advantage side & step-wise reward gain \(\Delta R_t\) & ODE readout 成本高，接近 \(O(T^2)\)\\
Chunk-GRPO & advantage side / 粒度 & chunk-level ratio 和 chunk-level update & 多数情况下仍依赖 terminal advantage\\
TempFlow-GRPO & advantage side / 探索 & branch-level counterfactual credit & branch 设计和额外采样成本复杂\\
PCPO & ratio side & 校准 Flow SDE ratio 的 hidden timestep weight & 不直接提供 step-wise reward gain\\
\bottomrule
\end{tabular}
}
\end{center}

这张表的重点不是“谁最好”，而是它们在拆同一个目标的不同部分：

\[
\textbf{让每个 timestep 的实际更新强度更接近它对最终质量的真实贡献。}
\]

\subsection{7. 当前形成的统一视角}

我们现在可以把 Flow-GRPO 的每个 step 更新写成更完整的概念式：

\[
\Delta \theta_t
\propto
\underbrace{
\nabla_\theta \log p_\theta(x_{t-1}|x_t,c)
}_{\text{policy gradient direction}}
\cdot
\underbrace{
A_t
}_{\text{advantage credit}}
\cdot
\underbrace{
w_{\rho}(t)
}_{\text{ratio / SDE scale}}
.
\]

标准 Flow-GRPO 的问题是：

\[
A_t
\approx
A_{\mathrm{terminal}},
\qquad
w_{\rho}(t)
\text{ 没有显式校准。}
\]

这就导致两类失真：

\begin{enumerate}[leftmargin=1.5em]
  \item \textbf{Advantage 侧失真：} 该 step 不一定真的贡献了 reward，但它拿到了同一个 \(A_i\)。
  \item \textbf{Ratio 侧失真：} 该 step 就算拿到了合理 \(A_t\)，它的 logprob / ratio scale 也可能让实际更新偏大或偏小。
\end{enumerate}

\begin{summarybox}
\textbf{最重要的阶段结论：}
\[
\boxed{
\text{Flow-GRPO 的 credit assignment = advantage side + ratio side。}
}
\]
DenseGRPO、Chunk-GRPO、TempFlow-GRPO 主要让我们看到 advantage 分布不平衡；PCPO 让我们看到 policy ratio 本身也会影响 credit。
\end{summarybox}

\subsection{8. 对 OTCA 的定位}

OTCA 可以暂时降级为旁支启发，而不是这条主线里的关键节点。

它有一个有用提醒：

\[
\text{多 reward objective 也会影响 advantage 怎么组织。}
\]

也就是除了问：

\[
\text{which timestep?}
\]

还可以问：

\[
\text{which objective?}
\]

但 OTCA 自己的实现没有前四篇那么硬：

\begin{itemize}[leftmargin=1.5em]
  \item TCD 的 latent similarity gain 只是 proxy，不是真正 causal reward gain。
  \item MOCA 主要是在 advantage scalar 上做融合，不是真正的 gradient-space MOO。
  \item 主公式里的 \(c_k^i\) 没有 \(t\)，所以“timestep-aware objective mixture”这个说法需要保留。
\end{itemize}

所以后续可以借 OTCA 的 when + what 框架，但不要直接照搬它的 TCD / MOCA。

\subsection{9. 后续读论文时的抓手}

之后继续读这条主线的论文，可以先问五个问题：

\begin{enumerate}[leftmargin=1.5em]
  \item 它是在修 \textbf{advantage side}，还是在修 \textbf{ratio side}？
  \item 如果修 advantage，它是 step、chunk、branch，还是 objective-level credit？
  \item 它的 credit 是 causal / counterfactual，还是 proxy / heuristic？
  \item 它有没有解决计算成本，还是只是把 credit 算得更密？
  \item 它是否考虑了 Flow SDE ratio 的 timestep scale，还是默认 ratio 是中性的？
\end{enumerate}

\begin{intubox}
\textbf{一个更高标准的目标：} 最理想的后续方法，应该同时做到：
\[
\begin{aligned}
&\text{低成本估计 temporal reward contribution}\\
&\quad +\ \text{校准 ratio-side timestep scale}\\
&\quad +\ \text{在多 reward 场景下避免 objective 冲突。}
\end{aligned}
\]
现在四篇论文各自解决了一部分，但还没有一篇把三件事都做干净。
\end{intubox}
